% Relatório do TCC2: edição mantida diretamente neste arquivo LaTeX.
\chapter{Fundamentação teórica}\label{cap:tcc2-2}

Este capítulo apresenta os conceitos necessários para caracterizar o produto desenvolvido. A discussão se concentra nos motores de jogos, nos jogos de tiro em primeira pessoa, nos elementos de plataforma e na organização de partidas multijogador competitivas em tempo real.

\section{Motores de jogos e Godot Engine}\label{sec:2.1}

Neste trabalho, o motor de jogos corresponde ao ambiente que reúne os recursos utilizados para construir e executar a aplicação. Na Godot, nós representam unidades de comportamento e podem ser organizados em cenas reutilizáveis. Uma cena de personagem, por exemplo, pode combinar corpo físico, câmera e apresentação visual. Essa composição também se aplica a menus e ao próprio mapa. A documentação da Godot apresenta nós, cenas, árvore de cenas e sinais como conceitos centrais de sua organização. \cite{GODOT-CONCEITOS}

Os comportamentos do jogo foram programados em GDScript. Para física tridimensional, entrada, interface, áudio e comunicação em rede, foram utilizados recursos do motor. A mesma base tecnológica é empregada nos dois projetos Godot, embora o cliente e o servidor possuam responsabilidades diferentes. O primeiro concentra a interação e a apresentação; o segundo mantém as regras e os dados compartilhados da partida.

\section{Jogos de tiro em primeira pessoa e arena shooters}\label{sec:2.2}

Nos jogos de tiro em primeira pessoa, a exploração e o combate são apresentados a partir do ponto de vista do personagem controlado. Hullett e Whitehead descrevem esse tipo de jogo como uma navegação por um mundo tridimensional em que o jogador enfrenta outros personagens utilizando armas. Essa perspectiva relaciona o enquadramento da câmera à direção de observação e de ataque. \cite{HULLETT}

O recorte de combate em arena adotado neste trabalho diz respeito às funcionalidades de movimentação, combate e coleta de recursos distribuídos pelo cenário. O jogo desenvolvido não inclui campanha narrativa ou progressão por níveis de personagem.

\section{Elementos de jogos de plataforma em ambientes 3D}\label{sec:2.3}

Os elementos de plataforma considerados no projeto são ações e estruturas que permitem vencer diferenças de altura e alcançar outras regiões do cenário. A proposta inicial incluiu salto duplo, plataformas suspensas e plataformas de impulso. A versão implementada preservou o salto duplo e as plataformas de impulso, relacionando o deslocamento vertical ao percurso pela arena.

\textit{Lovely Planet} oferece uma referência para a combinação entre saltar, mirar e disparar. A descrição do próprio desenvolvedor enfatiza a alternância entre desviar de projéteis e buscar uma posição de tiro. No presente trabalho, essa referência ajuda a contextualizar a intenção de reunir combate e movimentação, enquanto as regras concretas de salto e impulso são próprias da implementação. \cite{LOVELY}

\section{Jogos multijogador competitivos em tempo real}\label{sec:2.4}

O TCC1 organizou os jogos multijogador segundo a disposição dos participantes, sua forma de interação, os papéis desempenhados e o momento das ações. Dentro dessa classificação, o sistema desenvolvido pode ser descrito como um jogo em rede, competitivo, simétrico e de tempo real: os participantes executam instâncias em computadores conectados, disputam objetivos, utilizam o mesmo conjunto de ações e atuam durante a mesma partida. Nesse contexto, a simetria corresponde ao conjunto de ações e regras disponível a ambas as equipes.

A participação em tempo real não implica ausência de atraso. Cada computador possui seu próprio ciclo de execução, e as informações levam tempo para atravessar a rede. Por isso, a aplicação precisa conciliar a resposta percebida pelo jogador com o estado conhecido pelos demais participantes. Predição e interpolação são estratégias distintas para lidar com esse problema: a primeira antecipa localmente uma resposta; a segunda constrói uma transição entre estados recebidos. \cite{FIEDLER-REDE} \cite{FIEDLER-INTERPOLACAO}

